\begin{abstract}
A differentiable particle filter can compute the derivative of its own
finite likelihood accurately while estimating the underlying model score
poorly. This distinction explains both the progress and the negative results
of the BayesFilter investigation of Younis and Sudderth's kernel mixtures.
The original observation sidecar led to a complete kernelized-observation
recursion and then to an all-components, raw-importance-weight resampling
reference. The integrated derivatives pass their declared checks, but the
calibrated resampling construction has higher model-score mean squared error
than the canonical filter in both tested matrix-Gaussian scopes. We retain
these results and derive a different use of the same mixture ideas: construct
evaluable proposals, retain the original transition and observation densities
in the importance correction, and estimate the observed-data score by
averaging complete-data scores over possible ancestors. Selective pathwise
and importance-weight derivatives, valid control variates, and analytic
conditional integration supply additional variance-reduction candidates.
The note documents the implemented finite programs, the reasons for each
change in direction, the complete proposed recursions, and the experiments
needed to distinguish improved model-score accuracy from implementation
correctness. The new combination is a research proposal; its finite-particle
bias and computational value remain to be measured.
\end{abstract}

\section{The model score and the finite filter}
\label{sec:question}

The scientific question is how to estimate the score of a state-space model
accurately enough to support inference about its parameters. The existing
BayesFilter calculation uses \LEDH{} particle flow, PF--PF importance weights,
entropy-regularized \OT{}, Contract--E moment restoration, and \GenUT{}
higher-moment correction with two trust-region caps. These mechanisms define
a particular finite calculation. They are useful starting points for a
proposal, but they do not define what the exact model score must be.

This note grew out of the historical ``Section 3.6'' observation-sidecar
document. That earlier note remains in the isolated
\path{section-3-6-prior-weight-sidecar} worktree, under
\path{docs/papers/section_3_6_prior_weight_sidecar/}. The present manuscript,
\path{docs/papers/ledh_younis_kdm_score/ledh_younis_kdm_score.tex}, is its
broader successor. It brings together the subsequent implementation, results,
and the September 11 change in research question. The implementation history
is tied to research commit \texttt{804616e3}; later differences in the current
checkout, \texttt{5cc59cfa}, are identified in Section~\ref{sec:implementation}.

For fixed observations $y_{1:T}$ and model parameters $\theta\in\R^p$, write
\begin{align}
 \ell(\theta)&=\log p_\theta(y_{1:T}), &
 s(\theta)&=\nabla_\theta\ell(\theta), \label{eq:exact-score}\\
 L_0^N(\theta;\xi)&=\text{value of the executed finite filter}, &
 S_0^N(\theta;\xi)&=\nabla_\theta L_0^N(\theta;\xi).
 \label{eq:finite-score}
\end{align}
Here $N$ is the particle count and $\xi$ is the fixed random design used by
the finite program. In this note $\dd F$ means the directional derivative
$v^\mathsf{T}\nabla_\theta F$ for a declared direction $v$; $s$ denotes the
complete vector. Checking $S_0^N$ against finite differences of $L_0^N$ tests
the implementation of that derivative. Comparing it with $s$ tests its
accuracy as a model-score estimator. A third calculation, developed in
Section~\ref{sec:fisher}, estimates $s$ through posterior expectations. It
need not equal $S_0^N$ for the same particle realization.

For a fixed dataset, define the conditional bias and covariance of the
finite-program score by
\begin{equation}
 b_N(\theta;y)=\mathbb E_\xi[S_0^N\mid y]-s(\theta),\qquad
 V_N(\theta;y)=\operatorname{Cov}_\xi(S_0^N\mid y).
 \label{eq:conditional-score-error}
\end{equation}
When the second moments exist, its squared Euclidean error has expectation
$\|b_N\|^2+\operatorname{tr}V_N$. An individual discrepancy from Kalman does
not identify which term is responsible. Averaging errors over different
observation paths answers another question again: conditional biases can
cancel across datasets. Repeated filter randomizations within each dataset
are needed to distinguish conditional bias from Monte Carlo variance.

Even an unbiased likelihood estimate does not generally give an unbiased
log-likelihood score. Consider the positive estimator
$\widehat Z(\theta)=\theta+\varepsilon$, where
$\Pr(\varepsilon=a)=\Pr(\varepsilon=-a)=1/2$ and $\theta>a>0$.
Its mean is $Z(\theta)=\theta$, and its derivative is exactly one. Yet
\begin{equation}
 \mathbb E\left[\partial_\theta\log\widehat Z\right]
 =\frac12\left(\frac1{\theta+a}+\frac1{\theta-a}\right)
 =\frac{\theta}{\theta^2-a^2}\ne\frac1\theta
 =\partial_\theta\log Z.
 \label{eq:unbiased-value-biased-score}
\end{equation}
At $\theta=2$ and $a=1/2$, the two scores are $8/15$ and $1/2$.
The bias is created by the random denominator, despite correct
differentiation and an unbiased numerator. This example does not say that
every finite-particle score is biased in every model. It shows why the two
correctness statements cannot be substituted for a score-accuracy argument.

There are consequently several possible repairs. Missing derivatives of the
state, covariance, initial law, or sampling distribution require an
implementation or estimator correction. A reset or kernel that changes the
filtering measure requires a correction to that approximation or an explicit
new target. A valid but noisy integral requires better integration. The
development below separates these explanations before selecting the next
method. A zero-mean control variate addresses the last problem; it does not
remove an arbitrary bias already present in its baseline.

The existing implementation names remain useful for identifying what ran.
\texttt{ATOM-FINITE} denotes $L_0^N$; \texttt{KDM-FINITE} denotes a
positive-bandwidth observation program; \texttt{RESKDM-IWSG-FINITE} denotes
the post-Contract--E resampling program with raw anchored ratios; and
\texttt{RESKDM-SN-FINITE} denotes its historical, prematurely normalized
variant. \texttt{MODEL-IS} denotes an importance calculation whose numerator
retains the original model factors. These are program identities, not levels
of scientific validity. In particular, a correct \texttt{MODEL-IS} update
does not by itself make an additional deterministic reset unbiased, and a
different finite derivative is acceptable only under its actual name.

The \LEDH{} spelling follows the Li--Coates particle-flow construction
\citep{li2017particle}; the differentiable-transport context is
\citet{corenflos2021differentiable}. The new investigation now includes
alternative model-score estimators explicitly. It no longer rules them out
merely because they differentiate a different finite scalar.
